% This needs to come FIRST in the document.
\DocumentMetadata{
  lang        = en,
%	 pdfversion  = 1.7,
%  pdfstandard = ua-1,
  pdfstandard = ua-2,
  tagging=on,
%	 tagging-setup={math/mathml/luamml/load=false}
  tagging-setup={math/setup=mathml-SE}
}
\tagpdfsetup{table/header-rows={1}}

\documentclass[10pt]{article}

\usepackage[margin=1in]{geometry} % At least Geometry is legal with tagging!

\usepackage[base]{babel} % needed for lipsum in LuaLaTeX
\usepackage{lipsum} % for dummy text

\usepackage{tikz} % for graphics; the same commands work with \includegraphics

\title{PDF 2.0 Accessibility Demo}
\author{Masaros}
\date{April 10, 2026}

\begin{document}
\maketitle
This document provides an example of an accessible PDF using the PDF/UA-2 standard.
\section{First section}
\lipsum[1]
\section{A section with some math}
See how this version handles math formulas, like this displayed equation.
% Displayed math
\[
y=mx+b
\]
We can also write a paragraph with inline math, such as $y=mx+b$ or $x=\frac{2}{3}$.
\section{A section with some pictures and tables}
Let's draw an image with alt text.
\begin{center}
\begin{tikzpicture}[x=0.25in,y=0.25in,alt={A circle}]
\draw (0,0) circle (1);
\end{tikzpicture}
\end{center}
This next image is just here to look pretty, so I'm going to tell screen readers to ignore it.
\begin{center}
\begin{tikzpicture}[x=0.25in,y=0.25in,artifact] %"artifact" is a decorative image
\draw (-1,-1) rectangle (1,1);
\end{tikzpicture}
\end{center}
And finally, inline, I'm going to make an image that's just some text with a box around it.
\begin{tikzpicture}[baseline=0,actualtext=Text]
\draw (0,0) node[anchor=base,draw]{Text};
\end{tikzpicture}
The screen reader will ignore the box and just read the word.
\begin{center}
\begin{tabular}{lll}
Column 1&Column 2&Column 3\\\hline
Row 1&Row 1&Row 1\\
Row 2&Row 2&Row 2
\end{tabular}
\end{center}
\end{document}